\BeleskeID{09_1-s056}
% element: 09_1-s056:n01
Tranzistori na zajedničkoj osnovi ne mogu se posmatrati kao dva potpuno nezavisna tranzistora sa osnovom vezanom za sopstveni sors. Osnova je na masi, dok je sors opterećenja na izlazu. Zato se prag opterećenja menja sa izlaznim naponom usled efekta osnove. Jednačina za visoki izlazni nivo implicitna je: $V_{OH}$ se pojavljuje i izvan korena i u njemu. Može se rešiti svođenjem na kvadratnu jednačinu, uz proveru dobijenih korena, ili iterativno. Tačka $V_{IL}$ određena je pragom pogonskog tranzistora, pa u ovom razmatranju ne zavisi od praga opterećenja. Za $V_{IH}$ treba uzeti u obzir zavisnost praga opterećenja od napona sors–osnova, odnosno od izlaza.

\[V_{Tn}=V_{Tn0}+\gamma\left(\sqrt{|2\phi_F+V_{SB}|}-\sqrt{|2\phi_F|}\right)\]

\[V_{SB,L}=V_{S,L}-V_{B,L}=V_{OH}-0=V_{OH}\]

